% TOP_PROBLEM: {"id": 1289, "title": "Fermat-Catalan Conjecture", "category_id": 1, "category": {"id": 1, "name": "number_theory", "display_name": "Number Theory", "description": "Properties of integers, prime numbers, Diophantine equations.", "slug": "number-theory", "order_index": 1, "created_at": "2026-07-31T15:26:25.670Z"}, "status": "open"}
% FIRST_POSTED: 2026-09-25
% ENTRY_KIND: statement_only
% !TeX program = lualatex
% Definition and sources only; this file is not an LLM solution attempt.
\documentclass[11pt]{article}
\usepackage[margin=25mm]{geometry}
\usepackage{fontspec}
\setmainfont{Latin Modern Roman}
\usepackage{amsmath,amssymb,mathtools}
\usepackage{unicode-math}
\setmathfont{Latin Modern Math}
\usepackage{xurl,hyperref}
\hypersetup{colorlinks=true,urlcolor=blue}
\setcounter{secnumdepth}{0}
\setlength{\parindent}{0pt}
\setlength{\parskip}{5pt}
\setlength{\emergencystretch}{3em}
\begin{document}
\section*{\texorpdfstring{Fermat-Catalan Conjecture}{Fermat-Catalan Conjecture}}\label{1289}
\subsection{Definitions and mathematical statement}
For positive integers $x,y,z$ and integers $p,q,r>1$, consider
\[
 x^p+y^q=z^r,\qquad
 \gcd(x,y)=\gcd(y,z)=\gcd(z,x)=1,
 \qquad\frac1p+\frac1q+\frac1r<1.
\]
The recorded finiteness assertion concerns the set of base triples $(x,y,z)$ that admit at least one such exponent triple: this set is finite. The strict reciprocal-sum condition excludes the spherical and Euclidean exponent ranges. Ordering of $x,y$ is harmless for finiteness, but positivity and the coprimality restriction are essential. The exact source wording counts base triples; this must be distinguished from versions of Fermat--Catalan that count triples of powers $(x^p,y^q,z^r)$ or full sextuples. This entry does not silently identify those counting conventions.
\par\smallskip\textit{Formulation and concept references:} \href{https://home.cit.tum.de/~wvi/NT.pdf}{[S1]}.

\subsection{Short English statement}
Are there only finitely many coprime positive base triples whose powers satisfy an additive equation with reciprocal exponents summing to less than one?
\subsection{Sources}
\begin{itemize}
\item[S1]Violetta Weger, Elementary Number Theory, Conjecture 11.24 (2024/25).
\url{https://home.cit.tum.de/~wvi/NT.pdf}.
\textit{Formulation source.}
\end{itemize}
\end{document}
